\chapter{Limites et perspectives}

\section{Limites du modèle}

Le GRU de référence de la campagne expérimentale atteint 92,0 \% d'exactitude et un F1-score macro de 0,9181 sur les séquences du signeur réservé au test. Cette mesure concerne un seul découpage et ne garantit pas les mêmes performances pour d'autres signeurs. Le TFLite actuellement déployé dans l'application obtient séparément 90/100 sur les séquences prétraitées du jeu de test ; la correspondance entre ses poids et ceux du nouveau GRU de référence n'est pas établie. Avec quatre personnes dans le corpus, la variabilité des façons de produire les signes reste peu représentée.

Les résultats par classe montrent aussi que la performance varie selon le signe. Tsia et Maty font partie des signes les plus difficiles à reconnaître dans l'évaluation présentée au chapitre 3. Ces erreurs peuvent être liées aux différences de réalisation entre signeurs, mais le corpus ne permet pas d'établir une cause unique. Les points du corps et des mains employés par le modèle ne représentent pas toute l'information linguistique disponible, notamment les expressions du visage et les mouvements de la tête.

Le système traite des signes isolés associés à un vocabulaire fermé de dix classes. Il ne segmente pas une conversation en signes et ne traduit pas des phrases continues. Il ne prend donc pas en compte le contexte linguistique, l'ordre des signes ni les constructions propres à la LSM.

\section{Limites de l'application}

L'application est un prototype de démonstration et non un service déployé en production. La première utilisation de la synthèse vocale nécessite une connexion internet pour charger le modèle ONNX dans le navigateur. Les essais de reconnaissance dépendent également de la qualité de la vidéo et de la détection des landmarks par MediaPipe. Le cadrage, l'éclairage ou une main hors du champ peuvent réduire la quantité d'information disponible pour la prédiction.

La validation de l'application a porté sur un nombre limité de scénarios. Dans les dix essais de l'endpoint décrits au chapitre 4, l'audio est généré pour le résultat disponible, mais la vidéo du signe Miarahaba est classée Solosaina. Ce résultat rappelle que le fonctionnement de la chaîne logicielle ne garantit pas l'exactitude de chaque prédiction. Une campagne de tests plus vaste est nécessaire pour mesurer séparément les erreurs de capture, de détection, de classification et de synthèse.

\section{Limites méthodologiques}

Le corpus comprend 400 séquences pour dix signes et quatre signeurs. Il permet de réserver un signeur au test, mais sa taille ne couvre qu'une partie des variations possibles entre locuteurs. Les captures ont été effectuées dans deux pièces sous un éclairage naturel resté constant, mais le script ne contrôle pas ces conditions et ne documente pas systématiquement le fond ni la distance à la caméra. Les résultats ne permettent donc pas de conclure à la robustesse dans des lieux, éclairages et cadrages variés.

L'évaluation finale porte sur un seul signeur tenu à l'écart de l'entraînement et de la validation. Elle constitue un test de généralisation inter-signeurs, mais ne permet pas d'estimer la performance moyenne sur l'ensemble des personnes utilisant la LSM. Les résultats doivent être interprétés dans ce périmètre.

La campagne expérimentale repose sur un découpage fixe et sur un seul signeur réservé au test. Les écarts entre les trois meilleures exactitudes observées (94,0 \%, 93,0 \% et 92,0 \%) représentent au plus deux séquences sur les 100 du jeu de test. Aucun test statistique ni intervalle de confiance n'a été calculé ; ces résultats ne permettent donc pas d'établir une différence statistiquement significative ni une performance moyenne sur d'autres signeurs. Les comparaisons de taille de modèle sont également limitées par les taux de dropout différents entre configurations. Enfin, la variante à 45 frames est construite par rééchantillonnage des séquences déjà ramenées à 30 frames et n'ajoute pas d'information provenant des vidéos originales.

L'expérience d'augmentation combine du bruit gaussien et un décalage temporel circulaire. L'exactitude du GRU passe de 92,0 \% à 77,0 \%, tandis que celle du LSTM reste à 82,0 \%. Le F1-score macro du LSTM augmenté n'est pas disponible. L'effet propre de chaque transformation ne peut être isolé, et ce résultat ne permet pas de conclure à l'inefficacité de toute augmentation. L'analyse des séquences d'entraînement révèle une discontinuité potentielle aux extrémités dans 206 séquences sur 210 ; elle fournit une piste d'explication pour le jitter, sans établir de lien causal avec la baisse mesurée. Enfin, le TFLite actuellement déployé atteint 90/100 sur les caractéristiques de test, alors que le journal Colab du nouveau GRU de référence annonce 92/100. La correspondance entre les poids de ces deux artefacts n'est pas établie.

\section{Perspectives}

Une première perspective consiste à enrichir le corpus et à évaluer les modèles sur plusieurs signeurs tenus à l'écart, par exemple au moyen d'une validation où le signeur de test varie. Cette évaluation permettrait de mesurer la variabilité des résultats et de vérifier si les scores du modèle mains seules se maintiennent au-delà du découpage actuel. Les différences entre modèles pourraient alors être accompagnées d'intervalles d'incertitude et de comparaisons statistiques adaptées aux prédictions appariées.

Une deuxième perspective est d'étudier les variantes mains seules et à 15 frames comme pistes d'optimisation. Le modèle mains seules a le plus petit nombre de paramètres parmi les trois meilleures configurations, mais son avantage observé doit être confirmé sur un corpus plus large. La fenêtre de 15 frames pourrait réduire le nombre d'étapes récurrentes ; un benchmark sur la machine de démonstration reste nécessaire pour mesurer son effet réel sur la latence et la mémoire. Ces mesures permettraient également de comparer la taille et les performances des modèles après conversion en TFLite.

La reproductibilité du déploiement constitue une autre étape nécessaire. Il conviendrait d'archiver le checkpoint Colab correspondant au GRU de référence à 92,0 \%, de le convertir en TFLite, puis de vérifier que l'artefact obtenu correspond bien au modèle évalué. Le TFLite actuellement présent dans le dépôt obtient 90/100 sur les caractéristiques de test ; tant que la provenance des poids n'est pas reliée au nouveau checkpoint, ces deux scores doivent rester distincts.

Pour l'augmentation des données, des expériences d'ablation pourraient isoler le bruit gaussien du décalage temporel. Des transformations temporelles qui préservent l'ordre et la continuité du geste, ou des perturbations fondées sur la cinématique, pourraient ensuite être étudiées et évaluées sur des signeurs non utilisés pour le réglage. Les résultats actuels ne permettent pas de conclure à l'efficacité de ces méthodes.

À plus long terme, le système pourrait être étendu à la reconnaissance de signes continus. Cette évolution nécessiterait des séquences annotées avec des frontières temporelles et des traductions de référence. Des méthodes telles que la classification temporelle connexionniste (CTC) ou des architectures temporelles pourraient être comparées dans un protocole adapté. Une évaluation avec des utilisateurs de la communauté sourde malgache aiderait également à examiner la pertinence du vocabulaire, des traductions et de la sortie vocale.

\endinput
